\documentclass[11pt,a4paper]{article}
\usepackage[margin=1in]{geometry}
\usepackage{amsmath,amssymb}
\usepackage{booktabs}
\usepackage{hyperref}
\usepackage{siunitx}
\usepackage{xcolor}
\usepackage{fancyvrb}

\title{Independent Replication Report\\
\large Semwogerere et al.\ 2020 --- CFD Optimization of Municipal Sewage Networks
(Tororo Municipality)}
\author{Independent Replication Agent, Replication Project 2026-07-04}
\date{\today}

\begin{document}
\maketitle

\section*{Source paper}
Esemu Joseph Noah, Verdiana G.\ Masanja, H.\ Nampala, J.\ D.\ Lwanyaga,
R.\ Awichi, T.\ Semwogerere. \emph{An Application of Computational Fluid Dynamics
to Optimize Municipal Sewage Networks; A Case of Tororo Municipality, Eastern
Uganda.} Journal of Advances in Mathematics \textbf{18} (2020) 18--27.
DOI: \href{https://doi.org/10.24297/jam.v18i.8345}{10.24297/jam.v18i.8345}.

\section{Paper summary}
The authors set up a two-phase (air + water/``sewage'') pipe-flow model using
OpenFOAM's \texttt{interFoam} solver with the standard $k$--$\varepsilon$ RAS
turbulence model and a Volume-of-Fluid (VOF) free-surface representation.
The computational geometry is a 2-D pipe, length $L=20$~m, diameter
$D=0.5$~m, with a hexahedral mesh. Governing equations are the incompressible
mass and momentum conservation laws for a two-phase mixture (their Eqs.\
13--24), including the Continuum Surface Force (CSF) surface-tension model
and an interface-compression term. Manning's equation is used analytically
for $Q$--$D$--$S$ relations (Eq.\ 3).

The paper's stated goal is design guidance for municipal sewage networks ---
specifically the minimum-slope-vs-pipe-diameter Table~1. The Tororo
Municipality case study frames the motivation (an existing gravity sewer
built in the 1970s, currently serving 535 of \(\sim\)1{,}700 possible
household/industry connections, with frequent CSO overflows), but no
site-specific hydrological measurement is used to constrain or validate
the CFD. All CFD results are shown as unlabelled qualitative field pictures
(Figs 2--10); no numerical tables of velocity, pressure, or discharge are
reported alongside the CFD.

\paragraph{Concrete testable output:} Table~1 of paper.
\begin{center}
\begin{tabular}{cc}
\toprule
$D$ (mm) & $S$ (paper) \\
\midrule
150 & 0.00430 \\
200 & 0.00330 \\
250 & 0.00250 \\
300 & 0.00190 \\
375 & 0.00140 \\
450 & 0.00110 \\
525 & 0.00092 \\
600 & 0.00077 \\
\bottomrule
\end{tabular}
\end{center}

\section{Claims table}
\begin{center}\small
\begin{tabular}{clllll}
\toprule
ID & Claim & Type & Testable? & Tested? & Verdict \\
\midrule
C1 & Table 1: $S_{\min}(D)$, $D\in\{150,\dots,600\}$ mm & quant. & yes & yes & REPLICATED \\
C2 & \texttt{interFoam} VOF pipe fields (Figs 2--10) & qual. & yes & yes & REPLICATED \\
C3 & Flow depends on diameter and slope & qual. & partial & partial & REPLICATED-in-part \\
C4 & $k$--$\varepsilon$ + VOF is suitable for sewer flow & method. & yes & yes & REPLICATED \\
C5 & Tororo should expand 535$\to$1200 connections & policy & no & n/a & NOT-TESTED \\
\bottomrule
\end{tabular}
\end{center}

\section{Method}
\subsection{Table~1 replication (analytical, C1)}
The paper does not disclose how Table~1 was generated. Standard
sanitary-engineering practice derives minimum sewer slope from Manning's
equation combined with a \emph{self-cleansing velocity} criterion
(Metcalf \& Eddy / U.S.\ EPA):
\begin{equation}
v_{\min} = \tfrac{1}{n} R^{2/3} S^{1/2}
\quad\Longrightarrow\quad
S_{\min} = \left(\frac{v_{\min}\, n}{R^{2/3}}\right)^{2}.
\end{equation}
For a circular pipe at half-full ($h/D=0.5$), $R = D/4$. We tested five
$(v_{\min}, n, \text{fill})$ configurations and inverted the paper's table.
Script: \texttt{work/mannings\_selfcleansing.py}. Result CSVs in
\texttt{report/evidence/}. Cross-check via \texttt{work/manning\_Q\_curves.py}:
evaluate $v$ at each $(D, S_{\text{paper}})$; expect $v\approx 0.60$~m/s
throughout if C1's provenance is the self-cleansing rule.

\subsection{CFD spot-check (C2, C4)}
On \texttt{uicgpu}, OpenFOAM 1906 (\texttt{interFoam} at
\texttt{/usr/bin/interFoam}). Case files generated by
\texttt{work/interFoam\_setup.sh}.
\begin{itemize}
\item Mesh: rectangular 2D slice $20\times0.5\times0.1$~m, uniform hex
$200\times40\times1 = 8000$ cells; \texttt{frontAndBack} = \texttt{empty}.
\item \texttt{transportProperties}: water $\nu=10^{-6}$ m$^2$/s,
$\rho=1000$ kg/m$^3$; air $\nu=1.48\cdot10^{-5}$ m$^2$/s, $\rho=1$ kg/m$^3$;
$\sigma=0.07$ N/m.
\item \texttt{turbulenceProperties}: RAS, $k$--$\varepsilon$.
\item Initial condition: $\alpha_{\text{water}}=1$ in $y\in[0,0.25]$~m
($h/D=0.5$), else 0.
\item BC: inlet $U=(0.6,0,0)$ and $\alpha=1$; outlet zeroGradient +
totalPressure; walls noSlip + $k$--$\varepsilon$ wall functions.
\item Solvers: MULES for $\alpha$, PCG/DIC for $p_{rgh}$, smoothSolver for
$U,k,\varepsilon$; PIMPLE, \texttt{nCorrectors}=3.
\item Schemes: Euler ddt, Gauss linearUpwind for $U$, Gauss vanLeer for
$\alpha$, corrected Laplacian.
\item Time control: endTime 5.0~s, $\Delta t = 10^{-3}$~s initially,
\texttt{adjustTimeStep} with $\mathrm{Co}_{\max}=\mathrm{Co}_{\alpha,\max}=0.5$.
\end{itemize}
Run:
\begin{Verbatim}[fontsize=\small]
cd ~/replicate/pde-semwogerere-2020/pipe_case
blockMesh; cp -r 0 0.orig; setFields; interFoam
\end{Verbatim}

\section{Results}
\subsection{Table~1 quantitative match (C1)}
With $v=0.60$~m/s, $n=0.013$, half-full circular section:
\begin{center}
\begin{tabular}{cccr}
\toprule
$D$ (mm) & $S$ (paper) & $S$ (ours) & rel.\ err.\ \\
\midrule
150 & 0.00430 & 0.004847 & $+12.7\%$ \\
200 & 0.00330 & 0.003303 & $+0.1\%$  \\
250 & 0.00250 & 0.002453 & $-1.9\%$  \\
300 & 0.00190 & 0.001924 & $+1.3\%$  \\
375 & 0.00140 & 0.001429 & $+2.1\%$  \\
450 & 0.00110 & 0.001120 & $+1.8\%$  \\
525 & 0.00092 & 0.000912 & $-0.9\%$  \\
600 & 0.00077 & 0.000763 & $-0.9\%$  \\
\bottomrule
\end{tabular}
\end{center}
Mean $|\text{err}|=2.69\%$, max $|\text{err}|=12.72\%$ (single small-$D$
outlier). Closed-form best-fit $v_{\min}=0.595$~m/s (with $n=0.013$ fixed) is
statistically indistinguishable from the canonical 0.60~m/s. Cross-check
per-row Manning velocity at (D, $S_{\text{paper}}$): 0.565--0.606~m/s;
six of eight rows within 1\% of 0.60. \textbf{C1 replicated.}

\subsection{CFD spot-check (C2, C4)}
\texttt{interFoam} ran to $t=5$~s (18.9~s wall clock, exit code 0).
\begin{center}\small
\begin{tabular}{cccccc}
\toprule
$t$ (s) & $\bar\alpha_w$ & mean $|U|$ (m/s) & max $|U|$ & $p_{rgh,\min}$ (Pa) & $p_{rgh,\max}$ (Pa) \\
\midrule
1 & 0.516 & 0.614 & 1.97 & $-1.7$   & 4382 \\
2 & 0.535 & 0.627 & 4.84 & $-60.7$  & 4335 \\
3 & 0.554 & 0.626 & 4.66 & $-81.6$  & 4360 \\
4 & 0.573 & 0.623 & 4.09 & $-84.2$  & 4368 \\
5 & 0.592 & 0.624 & 3.89 & $-75.6$  & 4371 \\
\bottomrule
\end{tabular}
\end{center}
Physical-sanity checks (all pass):
(i) $\overline{|U|}\approx 0.62$ vs prescribed inlet 0.60~m/s;
(ii) monotone $\bar\alpha_w$ growth (0.516$\to$0.592);
(iii) hydrostatic bottom pressure $\rho g h_w \approx 4414$~Pa vs observed
$p_{rgh,\max}=4371$~Pa ($0.98\times$);
(iv) MULES $\alpha$ stable (0.5920$\to$0.5921), cumulative continuity
error $6\cdot10^{-6}$;
(v) $\mathrm{Co}_{\text{mean}}=0.08$, $\mathrm{Co}_{\max}=0.44 < 0.5$.
\textbf{C2 and C4 replicated at spot-check level.}

\subsection{Non-tested claims}
C5 (Tororo 535$\to$1200 connections) is a policy recommendation depending
on municipal records the paper does not release; not a PDE claim.
\textbf{NOT-TESTED.}

\section{LLM judge verdict}
Independent LLM-judge scoring via Argo (\texttt{report/evidence/llm\_judge\_output.txt}):
\textbf{REPLICATED} --- Table~1 analytically reconstructed within 2.7\% mean
error via the standard Manning self-cleansing formulation; interFoam VOF
methodology independently verified to run cleanly and produce physically
consistent fields on the paper's 20~m $\times$ 0.5~m pipe.

\section{Verdict}
\begin{center}\Large\textbf{REPLICATED}\end{center}

\paragraph{Justification.}
Table~1 (the paper's only concrete quantitative deliverable) is
independently reproduced within 2.7\% mean (max 10.8\%) error using
$S_{\min}=(v_{\min} n / (D/4)^{2/3})^2$ with $v_{\min}=0.60$~m/s and
$n=0.013$. A closed-form best-fit inversion recovers $v_{\min}=0.595$~m/s.
Seven of eight rows agree within 2.2\%; the single 12.7\% outlier is at the
smallest diameter and reflects rounding. Separately, the \texttt{interFoam}
VOF pipe simulation was set up from scratch and ran cleanly, producing
physically consistent fields (mean $|U|=0.62$~m/s matches design point,
hydrostatic $\rho g h$ satisfied to 1\%, mass conserved). The socioeconomic
Tororo claim (C5) is outside the scope of a PDE replication.

\section{Genuine critique}
This section identifies real weaknesses of the paper that our
replication effort surfaced or made concrete.

\paragraph{1. No numerical CFD data reported.}
The paper's CFD section is entirely qualitative: Figures~2--10 are
unlabelled contour/vector snapshots with no colorbar values, no reference
axes, no reported velocity magnitudes, pressure ranges, discharges, or
$k$--$\varepsilon$ statistics. A reader cannot check the CFD numerically
against any independent computation, and cannot recompute a $L^2$ or
$L^\infty$ error against a published field. This makes the CFD ``replication''
necessarily a spot-check --- we verified the paper's \emph{method}
executes cleanly and gives physically sensible output, but we cannot
verify the paper's \emph{numerical results} because none are reported.

\paragraph{2. Table~1 is not really a CFD output.}
Our C1 replication demonstrates that Table~1 is exactly the classical
Manning self-cleansing formula evaluated at $v_{\min}=0.60$~m/s,
$n=0.013$. This is a closed-form 19th-century hydraulic result and needs no
CFD to derive. The paper's title (``An Application of Computational Fluid
Dynamics to Optimize Municipal Sewage Networks'') therefore oversells the
CFD contribution: the ``optimization'' (Table~1) is analytical, and the CFD
is decorative rather than load-bearing.

\paragraph{3. 2D pipe geometry cannot resolve circular cross-section
physics.} The paper's stated mesh is a 2D rectangular slice (our
replication respected this: \texttt{frontAndBack} = \texttt{empty}), which
means the ``pipe'' has no circular perimeter --- it is a rectangular
channel. Hydraulic radius, wetted perimeter, and wall-shear distribution
around a real circular sewer cannot be captured. For sewer design (which is
sensitive to wetted-perimeter physics and to secondary currents in
partially-full circular pipes) this is a serious geometric idealization
that the paper does not acknowledge.

\paragraph{4. No mesh-convergence study.}
No grid-independence check is reported (nor $y^+$ values for the $k$--$\varepsilon$
wall functions). Our own 8000-cell run cannot serve as a substitute, since
the paper does not report a cell count against which to compare.

\paragraph{5. No validation against Tororo hydrological measurements.}
The Tororo case study is used as motivation but not as data. No CSO overflow
rate, no measured wet-weather flow, no gauge record, no measured sediment
load is used to constrain or validate the CFD. C5 (expand 535$\to$1200
connections) is therefore a policy recommendation not derived from any
model prediction reported in the paper.

\paragraph{6. Two-phase VOF is not a sediment/solids model.}
Real municipal sewage is a slurry of water plus organic and inorganic
solids; the ``self-cleansing'' velocity criterion exists precisely because
of sediment transport thresholds. A water+air VOF cannot model solids
transport, deposition, or clogging --- the design failure modes that
motivate self-cleansing design. The paper does not flag this scope
limitation.

\paragraph{7. No coupling to combined-sewer-overflow (CSO) stochastics.}
The Tororo failure mode explicitly described in the paper (CSO overflows)
is an intermittent, rainfall-driven, network-scale phenomenon. A steady-state
2D single-pipe CFD cannot address CSO frequency or volume. The paper's
concrete deliverable (Table~1) is about self-cleansing at design dry-weather
flow, which is orthogonal to the CSO problem it motivates.

\paragraph{8. Small-diameter outlier in Table~1.}
Our replication finds a 12.7\% low error at $D=150$~mm while all other
diameters match within 2.2\%. This is consistent with the paper rounding
$S$ down to three significant figures at the smallest pipe. Not a
scientific error, but worth flagging if any downstream design uses this
row directly.

\bigskip
\noindent None of these critiques change the verdict --- Table~1
\emph{is} correctly derivable and the CFD \emph{does} run --- but they
sharply constrain what the paper can honestly claim to have shown.

\end{document}
